(1) 周期边界使每个波数态的密度为 $L/(2\pi)$，但同一正频率 $\omega$ 对应 $k=\pm\omega/v$ 两支波数，因此每长度的正频率态密度为 $g(\omega)/L=1/(\pi v)$。对 $\omega>0$ 积分时不可再次任意乘以二。
(2) 由 Bose 平均占据 $\bar n=(e^{\beta\hbar\omega}-1)^{-1}$，热声子内能每长度为 $U/L=\frac1{\pi v}\int_0^\infty[\hbar\omega/(e^{\beta\hbar\omega}-1)]d\omega$。令 $x=\hbar\omega/(k_BT)$，得到 $U/L=(k_BT)^2/(\pi\hbar v)\cdot(\pi^2/6)=\pi k_B^2T^2/(6\hbar v)$。
(3) 对定长求导得 $C/L=\partial_T(U/L)=\pi k_B^2T/(3\hbar v)$。因为 $dS=(C/T)dT$ 且取 $S(0)=0$，有 $S/L=\pi k_B^2T/(3\hbar v)$；两者均随 $T\to0$ 消失，满足不考虑基态简并时的第三定律。
(4) 一维线性色散的态密度随频率为常数，故内能标度是 $T^2$、热容标度是 $T$；三维声学模态密度随 $\omega^2$ 增长，才导致 $C\propto T^3$。要求热能高于有限长最低模间距以便积分近似，同时 $k_BT\ll\hbar\omega_D$ 使高频色散截止可忽略。
